% ============================================================
% Chapter 5 -- Beyond Gaussianity: Laplace innovations
% ============================================================

\chapter{Beyond Gaussianity: Laplace Innovations and the Gaussian-Message
Approximation}\label{ch:laplace}\label{sec:bp-nongaussian}
% The full title wraps the running head onto a second line; the mark carries
% a one-line form.
\chaptermark{Beyond Gaussianity: Laplace Innovations}

Every exact result of Chapters~\ref{ch:gaussian} and~\ref{ch:gaussian-bp}
rests on the closure of the Gaussian family under the two operations the
model performs, multiplying densities and marginalising. This chapter removes
that property in the most controlled way available and answers RQ3: the AR(1)
recursion, the OU channel and the factor-graph structure are kept, and only
the innovation law changes, from Gaussian to Laplace.

\section{What changes, and what is preserved}\label{sec:l-model}\label{sec:l-nonlinear}

The Laplace chain keeps the initial law and the second-order structure of
\eqref{eq:ar1} fixed,
\begin{equation}
  a_0 \sim \Nd(0,1), \qquad
  a_k = \alpha\, a_{k-1} + \varepsilon_k, \qquad
  \varepsilon_k \sim \mathrm{Laplace}(0, b) \ \text{i.i.d.}, \qquad
  b = \sqrt{(1-\alpha^2)/2},
  \label{eq:l-chain-prior}
\end{equation}
so that $\operatorname{Var}(\varepsilon_k) = 1-\alpha^2$, the chain has unit
marginal variance and $\Cov(a_j,a_k) = \alpha^{|j-k|}$ exactly as before. The
Gaussian and Laplace chains agree on every second moment and differ only
beyond them. The score--posterior identity holds for any prior, so with
$\mu := e^{-t}$,
\begin{equation}
  \score_k(x,t) \;=\; \frac{\mu\, \E[a_k \mid x] \;-\; x_k}{\Deltat} ;
  \label{eq:l-score-identity}
\end{equation}
the posterior still factorises over the tree of
Figure~\ref{fig:posterior-factor-graph}, so the sweeps
\eqref{eq:bp-fwd}--\eqref{eq:bp-bwd} still compute exact smoothing marginals,
and the OU likelihood factors are still Gaussian in $a_k$. What is lost is
that the messages are no longer Gaussian and the posterior mean, hence the
score, is no longer linear in $x$ (Figure~\ref{fig:l-nonlinear}).

The two chains first differ at fourth order. A centred Laplace variable of
variance $v$ has fourth cumulant $3v^2$, excess kurtosis $3$; the chain
\emph{state}, a weighted sum of one Gaussian seed and $k$ Laplace innovations,
inherits a milder and position-dependent value, since fourth cumulants of
independent sums add with the fourth power of the coefficients:
\begin{equation}
    \kappa_4(a_k)
    = 3(1-\alpha^2)^2\sum_{m=0}^{k-1}\alpha^{4m}
    = 3\,\frac{1-\alpha^2}{1+\alpha^2}\,\bigl(1-\alpha^{4k}\bigr),
    \qquad
    \boxed{\;\kappa_4(x_{t,k}) = e^{-4t}\,\kappa_4(a_k)\;}
    \label{eq:l-state-kurtosis}
\end{equation}
the second identity because the channel adds an independent Gaussian with
$\kappa_4 = 0$. At $\alpha = 0.85$ the bulk value is $\approx 0.48$, not $3$;
the chain is covariance-stationary by construction but not strictly
stationary, because the Gaussian initial law is not the stationary law of the
Laplace recursion; and Laplace-driven non-Gaussianity survives every finite
diffusion time but decays as $e^{-4t}$, while the second moment is preserved
exactly. Both facts return in Chapters~\ref{ch:em-kernel}
and~\ref{ch:em-results}.

\begin{figure}[t]\centering
\includegraphics[width=\linewidth]{fig_laplace_nonlinear.pdf}
\caption[Laplace innovations: same variance, nonlinear score]{Laplace
(vermilion, solid) against Gaussian (blue, dashed) at $\alpha = 0.85$.
\emph{Left:} the two innovation laws on a log scale, same variance
$1-\alpha^2$, different shape. \emph{Centre and right:} the joint score
\eqref{eq:l-score-identity} at the middle frame of a chain of $\nsites = 9$
as $x_k$ is swept with the other frames at zero, by grid BP
(Section~\ref{sec:l-grid}) against the covariance-matched Gaussian chain,
which is exactly what the Gaussian closure returns
(Theorem~\ref{thm:l-lmmse}). At $t = 0.05$ the Laplace score saturates near
$\pm 5.2$ while the Gaussian one grows without bound; at $t = 1$ the two
coincide on the same scale.}
\label{fig:l-nonlinear}
\end{figure}

\section{Where closed forms stop}\label{sec:l-chain}

\begin{theorem}[Joint score from chain messages]\label{thm:l-markov-score}
For any chain prior, in particular \eqref{eq:l-chain-prior}, and the OU
likelihood, the joint score is \eqref{eq:l-score-identity} with
\begin{equation}
  \E[a_k \mid x]
  = \frac{\displaystyle\int a_k\; \lmsg_k(a_k)\;
      \psi^{(k)}_{\mathrm{ob}}(a_k)\; \rmsg_k(a_k)\;\dd a_k}
     {\displaystyle\int \lmsg_k(a_k)\;
      \psi^{(k)}_{\mathrm{ob}}(a_k)\; \rmsg_k(a_k)\;\dd a_k},
  \label{eq:l-markov-score}
\end{equation}
the messages being those of Proposition~\ref{prop:bp-convA}.
\end{theorem}

\noindent Computational locality survives; what does not survive is the
collapse of a message to two parameters. The first message can still be
written down: with $M(a'\mid a) = (2b)^{-1}e^{-|a'-\alpha a|/b}$ and
$y = x' - \mu\alpha a$, splitting the convolution
$\int M(a'\mid a)\,\Nd(x'; \mu a', \Deltat)\,\dd a'$ at the single kink and
completing the square on each half-line gives
\begin{equation}
  I(y) = \frac{e^{\Deltat/(2\mu^2b^2)}}{2\mu b}
  \Big[ e^{-y/(\mu b)}\,\Phi\!\Big(\tfrac{y}{\sqrt{\Deltat}} - \tfrac{\sqrt{\Deltat}}{\mu b}\Big)
      + e^{y/(\mu b)}\,\Phi\!\Big(-\tfrac{y}{\sqrt{\Deltat}} - \tfrac{\sqrt{\Deltat}}{\mu b}\Big) \Big],
  \label{eq:l-first-message}
\end{equation}
with $\Phi$ the standard Gaussian cdf. Closure is lost at the very next
update: the incoming message is now a sum of terms $e^{ca}\Phi(da+e)$, and
integrating it against the next Laplace transition requires half-line
integrals of $e^{-qa^2+\ell a}\Phi(da+e)$, higher-dimensional Gaussian
probabilities outside the normal--Laplace family, so each step increases the
functional complexity; changing to innovation coordinates only moves the
dependence from the prior into a dense likelihood. Exact BP remains valid,
and its messages must be represented numerically.

\section{Grid belief propagation}
\label{sec:l-grid}

A message is a density on $\R$. It is represented by its values on $\ngrid$
equispaced points $g_i = -A + ih$, $h = 2A/(\ngrid-1)$, of a truncated
interval $[-A, A]$, and every integral of
\eqref{eq:bp-fwd}--\eqref{eq:bp-bwd} becomes a composite trapezoidal
quadrature with weights $w_i = h$ in the interior and $h/2$ at the two
endpoints. The transition enters as a matrix and the likelihood as a diagonal,
$K_{ij} = M(g_i \mid g_j)$ and $(\ell_k)_i = \Nd(x_k; \mu g_i, \Deltat)$, and
with $\odot$ the elementwise product the sweeps read
\begin{align}
  f_0 &\;\propto\; p_0 \odot \ell_0,
  &f_k &\;\propto\; \ell_k \odot K\big( w \odot f_{k-1} \big),
  \quad k = 1,\dots,\nsites-1,
  \label{eq:l-grid-fwd}\\[2pt]
  \beta_{\nsites-1} &\;=\; \mathbf{1},
  &\beta_k &\;\propto\; K^\T \big( w \odot \ell_{k+1} \odot \beta_{k+1} \big),
  \quad k = \nsites-2,\dots,0 ,
  \label{eq:l-grid-bwd}
\end{align}
with belief $\mathrm{bel}_k = f_k \odot \beta_k$, posterior mean
$\widehat m_k = \sum_i w_i g_i (\mathrm{bel}_k)_i / \sum_i w_i (\mathrm{bel}_k)_i$
and score $\widehat\score_k = (\mu\widehat m_k - x_k)/\Deltat$. Each step is
one matrix--vector product, so the solver costs $O(\nsites\ngrid^2)$ time and
$O(\nsites\ngrid)$ space: sum--product over a finite alphabet, the alphabet
now a discretisation. Three choices fix the representation
(Table~\ref{tab:l-grid-choices}).

\begin{center}
\begin{minipage}{\linewidth}\small\centering
\captionof{table}[The three choices of the grid solver]{The three choices
that fix grid BP, what each is set to, and what fixes it.}
\label{tab:l-grid-choices}
\begin{tabular}{@{}p{0.2\linewidth}p{0.24\linewidth}p{0.5\linewidth}@{}}
\toprule
choice & setting & what fixes it \\
\midrule
renormalisation & forward messages to unit mass, backward to unit maximum, after every step
  & free, since messages are defined up to a constant (Theorem~\ref{thm:tree-exact}) and the ratio $\widehat m_k$ is invariant; necessary, since without it a chain of $\nsites = 32$ underflows \\
truncation $A$ & $A = 8$
  & the exponential tail beyond $|a| = 8$ carries below $10^{-9}$ of the mass for every $\alpha$ used; the mass in the outermost cells is logged per run \\
resolution $\ngrid$ & $\ngrid = 401$ (reference $801$)
  & the likelihood in $a$-coordinates has width $\sqrt{\Deltat}/\mu \approx \sqrt{2t}$ at small $t$; the gate
  \begin{equation}
    \sqrt{2t} \;\gtrsim\; 3h = 6A/(\ngrid - 1)
    \label{eq:l-resolution}
  \end{equation}
  keeps it wider than the spacing, and it is this gate, not the innovation law, that sets the smallest $t$ reported \\
\bottomrule
\end{tabular}
\end{minipage}
\end{center}

\paragraph{Validation.} Grid BP is a numerical approximation of the
functional recursion and serves as a reference only inside a validated
regime, established on the Gaussian chain where Chapter~\ref{ch:gaussian}
gives the truth. For smooth, rapidly decaying integrands the trapezoidal rule
converges superalgebraically, so failure is abrupt, and the measurement shows
exactly that: at $\ngrid = 101$, $A = 8$ the relative score error sits at the
double-precision floor down to $t \approx 0.03$, is $10^{-10}$ at $t = 0.02$
and reaches the percent scale by $t = 0.005$. The gate
\eqref{eq:l-resolution} is conservative relative to that onset, and the
working $\ngrid = 401$ clears it with room to spare ($3h = 0.12$ against
$\sqrt{2t} = 0.20$ at the smallest reported $t = 0.02$). For the Laplace
kernel the integrand has a kink and no spectral statement is invoked; the
reference is certified by self-convergence, the error between $\ngrid = 401$
and $801$ staying three to five orders of magnitude smaller than every effect
measured below (Figure~\ref{fig:bp-closure}, left). Self-convergence cannot
test the domain width, which the half-width sweep of \codefile{exp\_18}
($\gdtrunccells$ cells over $A \in \{4,6,8,10\}$) settles: the posterior mass
reaching the outermost cells is $\gdtruncfour$ at $A = 4$, $\gdtruncsix$ at
$A = 6$, $\gdtrunceight$ at the working $A = 8$ and $\gdtruncten$ at
$A = 10$, two orders below the smallest closure error reported here. Within
this regime grid BP is treated as the reference score; outside it no claim is
made.

One identity governs how everything is reported. For any two solvers with
posterior means $\widehat m$ and $m_{\mathrm{ref}}$,
\eqref{eq:l-score-identity} gives exactly
\begin{equation}
  \widehat s - s_{\mathrm{ref}}
  = \frac{e^{-t}}{\Deltat}\,\big( \widehat m - m_{\mathrm{ref}} \big),
  \qquad
  \frac{e^{-t}}{\Deltat} \sim \frac{1}{2t} \;\; (t \downarrow 0):
  \label{eq:bp-master-error}
\end{equation}
every score error is an amplified posterior-mean error, its residual logged in
every experiment between $10^{-15}$ and $10^{-12}$, and small diffusion times
are the difficult regime because whatever an approximation does to the
posterior mean is amplified by $1/2t$ in the score.

\section{The Gaussian-message approximation, measured}
\label{sec:l-closure}

\emph{Gaussian-projected BP} runs the same forward--backward recursion as
grid BP but projects every transition message onto the Gaussian family by
moment matching: $\Pi_{\mathrm G} f$ is the Gaussian with the mean and
variance of $f/\!\int f$. The first result says what that computes.

\begin{theorem}[Gaussian closure is the LMMSE estimator]
\label{thm:l-lmmse}
Consider moment-projected sum--product on the chain
\eqref{eq:l-chain-prior} with $\alpha \ne 0$, Gaussian unary likelihoods,
$\E[a_0] = 0$, innovations of mean zero and variance $\ivar$, and a Gaussian
(or projected) initial message. Its posterior means coincide exactly with
those of the covariance-matched Gaussian model,
\begin{equation}
  \widehat a_{\mathrm{LMMSE}}
  = \Cov(a,x)\Cov(x)^{-1}x
  = e^{-t}\Sigz \bigl(e^{-2t}\Sigz+\Deltat I\bigr)^{-1}x .
  \label{eq:l-lmmse}
\end{equation}
Hence the closure returns the same estimator for every innovation law sharing
the moments $(0,\ivar)$.
\end{theorem}

\begin{proof}
By induction every projected message equals the corresponding exact message
of the covariance-matched Gaussian chain. The initial messages coincide by
hypothesis. Multiplication by the Gaussian likelihood is the same
precision-weighted update in both models. A Gaussian message with moments
$(m, v)$ pushed through $a_i = \alpha a_{i-1} + \varepsilon_i$ has moments
$(\alpha m,\ \alpha^2 v + \ivar)$, and the backward transition integral,
viewed as a density in the preceding state, has moments
$(m/\alpha,\ (v + \ivar)/\alpha^2)$; both maps depend on the innovation law
only through $(0, \ivar)$, so after $\Pi_{\mathrm G}$ the projected message is
exactly the Gaussian chain's. (The terminal message $\rmsg_{\nsites-1} \equiv 1$
is not normalisable; the backward induction starts at the terminal likelihood,
the first object actually propagated, and $\alpha \ne 0$ is used in the
backward map.) All beliefs therefore coincide with the Gaussian chain's, whose
posterior mean is $\Cov(a,x)\Cov(x)^{-1}x$ with $\Cov(a,x) = e^{-t}\Sigz$ and
$\Cov(x) = e^{-2t}\Sigz + \Deltat I$.
\end{proof}

Gaussian-projected BP is therefore not a heuristic Gaussianisation but the
best second-order linear estimator, and comparing it with grid BP measures
the loss from discarding structure beyond the first two moments. The
\emph{estimator} is the same for every innovation law with matched moments;
its \emph{error} against the true posterior is not, because each law has its
own nonlinear posterior, and Section~\ref{sec:l-innovation-sweep} measures how
differently the same estimator misses different truths.

\paragraph{Three metrics.} The relative score error is what the diffusion
model consumes, but by \eqref{eq:bp-master-error} it is the posterior-mean
error multiplied by $e^{-t}/\Deltat$ and divided by the norm of the reference
score, both of which move by orders of magnitude across the schedule, so on
its own it cannot distinguish an approximation that improves from one whose
errors are scaled away. The primary metric is therefore the posterior mean,
which carries a risk interpretation.

\begin{proposition}[Excess-risk decomposition]\label{prop:l-excess}
For $m(x) = \E[a\mid x]$ and any estimator $\widehat a(x)$ with finite second
moment,
$\E\|a - \widehat a(x)\|^2 = \E\|a - m(x)\|^2 + \E\|m(x) - \widehat a(x)\|^2$.
\end{proposition}

\begin{proof}
Expand $a - \widehat a = (a-m) + (m - \widehat a)$; the cross term vanishes
after conditioning on $x$, since $\E[a - m(x)\mid x] = 0$ and both $m$ and
$\widehat a$ are functions of $x$.
\end{proof}

\noindent With $\widehat a = \widehat a_{\mathrm{LMMSE}}$ the gap between the
two posterior means is exactly the excess denoising risk of using second
moments only. It is reported normalised by the signal energy
$\E\|a\|^2 = \nsites$,
\begin{equation}
  \varepsilon_{\mathrm{PM}}
  \;=\;
  \sqrt{\tfrac{1}{\nsites}\,
        \E\big\| \widehat a_{\mathrm{LMMSE}} - \E[a \mid x] \big\|^2 } ,
  \label{eq:l-epspm}
\end{equation}
dimensionless, free of the prefactor and stable at both ends of the schedule.
Two secondary metrics accompany it: the excess-risk fraction
$\E\|m - \widehat a\|^2/\E\|a - m\|^2$, the percentage by which closure raises
the denoising error above the floor set by the posterior itself, and the
relative score error, because it is what the generative process sees.

\begin{center}
\begin{minipage}{\linewidth}\centering
\captionof{table}[Gaussian-closure error across diffusion time]{Gaussian
closure against validated grid BP on the Laplace chain ($\alpha=0.85$,
$\nsites=32$, $60$ trials per $t$, reference $\ngrid=801$), one sweep
sampled densely for $t\lesssim0.3$. Left group: the posterior approximation
before the prefactor, $\varepsilon_{\mathrm{PM}}$ \eqref{eq:l-epspm} with
its $10$--$90\%$ band, the same normalised to $t=0.02$, and the excess MSE of
Proposition~\ref{prop:l-excess} as a percentage of the MMSE floor. Last
column: the resulting discrepancy in the score.}
\label{tab:l-closure}
\small
\setlength{\tabcolsep}{5pt}
\renewcommand{\arraystretch}{1.1}
\begin{tabular}{@{}lccccc c@{}}
\toprule
& \multicolumn{5}{c}{\textit{posterior inference}}
& \textit{score} \\
\cmidrule(lr){2-6}\cmidrule(l){7-7}
$t$
& $\varepsilon_{\mathrm{PM}}$
& $10$--$90\%$
& $\varepsilon_{\mathrm{PM}}/\varepsilon_{\mathrm{PM}}^{\,0.02}$
& \shortstack{excess MSE\\{}[\% floor]}
& \shortstack{MMSE\\floor}
& \shortstack{rel.\ score\\error} \\
\midrule
$0.02$  & $0.036$  & $(0.029,\,0.052)$   & $1.00$ & $3.9$   & $0.032$ & $0.38$ \\
$0.032$ & $0.040$  & $(0.034,\,0.057)$   & $1.14$ & $3.4$   & $0.047$ & $0.30$ \\
$0.05$  & $0.050$  & $(0.036,\,0.079)$   & $1.42$ & $3.9$   & $0.066$ & $0.27$ \\
$0.072$ & $0.053$  & $(0.038,\,0.081)$   & $\mathbf{1.50}$ & $3.1$ & $0.089$ & $0.20$ \\
$0.105$ & $0.050$  & $(0.037,\,0.100)$   & $1.41$ & $2.1$   & $0.114$ & $0.14$ \\
$0.153$ & $0.043$  & $(0.034,\,0.085)$   & $1.20$ & $1.2$   & $0.151$ & $0.10$ \\
$0.222$ & $0.038$  & $(0.033,\,0.058)$   & $1.06$ & $0.74$  & $0.193$ & $0.065$ \\
$0.322$ & $0.033$  & $(0.027,\,0.059)$   & $0.93$ & $0.43$  & $0.253$ & $0.041$ \\
\addlinespace
$0.98$  & $0.016$  & $(0.012,\,0.028)$   & $0.45$ & $0.045$ & $0.584$ & $7.3\times10^{-3}$ \\
$2.07$  & $0.0018$ & $(0.0013,\,0.0038)$ & $0.05$ & $<0.01$ & $0.916$ & $2.4\times10^{-4}$ \\
\bottomrule
\end{tabular}
\end{minipage}
\end{center}

\paragraph{The measurement.} Table~\ref{tab:l-closure} and
Figures~\ref{fig:l-metrics} and~\ref{fig:bp-closure} report the sweep, and
two readings of it disagree. Read the last column and the closure improves
steadily as noise is added: the relative score error falls by a factor of
$3.6$ between $t = 0.02$ and $t = 0.15$. Read $\varepsilon_{\mathrm{PM}}$ and
it does not: the posterior-mean error \emph{rises} by half over
$t \in [0.02, 0.07]$, is still larger at $t = 0.15$ than at $t = 0.02$, and
begins to decay only past $t \simeq 0.3$. The gap between the readings is the
prefactor of \eqref{eq:bp-master-error} and the shrinking norm of the
reference score, neither of which reports on the inference. As a denoiser,
the excess over the MMSE floor falls from $3.9\%$ at $t = 0.05$ to $0.43\%$
at $t = 0.32$ and below $0.05\%$ once $t \gtrsim 1$, while the floor itself
grows from $0.07$ to $0.92$: the approximation error shrinks as the
irreducible uncertainty grows. Direction survives magnitude: the probability
that a trial exceeds $10\%$ relative score error is $1.00$ through
$t \simeq 0.07$ and zero by $t \simeq 0.7$, whereas the median cosine
similarity to the reference score is already $0.964$ at $t = 0.05$ and above
$0.99$ by $t \simeq 0.11$. Gaussian closure loses how strongly to move before
it loses where to move. The peak in $\varepsilon_{\mathrm{PM}}$ separates two
easy regimes: as $t \to 0$ each latent state is pinned by its own
near-noiseless observation, and at large $t$ the Gaussian noise washes out
the higher-order structure of the innovation law; the largest error sits
where the evidence is weak enough for the prior's shape to matter and strong
enough not to have been erased. The Laplace posterior is log-concave and
unimodal at every $t > 0$, so its worst trials are heavy-tailed and kinked
rather than bimodal; this is a qualitative mechanism for the aggregate curve,
not a measurement of posterior non-Gaussianity.

\begin{figure}[t]
\centering
\includegraphics[width=\linewidth]{fig_laplace_metrics.pdf}
\caption[The Gaussian closure under three metrics]{The same experiment under
three metrics ($\alpha=0.85$, $\nsites=32$, $60$ trials per $t$).
\emph{Left:} the posterior-mean error \eqref{eq:l-epspm}, median and
$10$--$90\%$ band, not monotone: it peaks at $t\simeq0.07$.
\emph{Centre:} the same runs through the relative score error, which falls
by $3.6\times$ over the interval on which $\varepsilon_{\mathrm{PM}}$ is
flat. \emph{Right:} the excess denoising MSE of
Proposition~\ref{prop:l-excess} as a percentage of the MMSE floor.}
\label{fig:l-metrics}
\end{figure}

\begin{figure}[t]
\centering
\includegraphics[width=\linewidth]{fig_laplace_closure.pdf}
\caption[Gaussian-closure score error against grid BP]{The score-side view.
\emph{Left:} median relative score error with its $10$--$90\%$ band, against
the grid-BP self-convergence error ($\ngrid=401$ versus $801$) on identical
data; the closure error exceeds the numerical budget by three to five orders
of magnitude, so the measurement is limited by the closure, not the
quadrature. \emph{Right:} relative score error against $1-\cos$ of the angle
between projected and reference score vectors; directional error is smaller
than magnitude error at every $t$, by a factor of about five at $t=0.02$ and
by three to five orders of magnitude once $t\gtrsim1$.}
\label{fig:bp-closure}
\end{figure}

\paragraph{Varying the innovation law.}\label{sec:l-innovation-sweep}
Holding $\Cov(a_i,a_j)=\alpha^{|i-j|}$ fixed, the innovation law is varied
over Laplace, Student-$t$ with $\nu\in\{5,8\}$, and symmetric Gaussian
mixtures of increasing separation. Across every family substantial closure
error is confined to $t\lesssim0.4$ and falls to about $10^{-3}$ or below by
$t=2.4$. At $t=0.05$ and $\alpha=0.5$ the median relative score error is
$0.94$ for a strongly bimodal mixture, $0.52$ for Laplace, $0.32$ for
Student-$t$ with $\nu = 5$, and $0.07$ for a mildly bimodal mixture. Strong
bimodality is the hardest case, a single Gaussian sitting between two
separated posterior modes. Stronger chain correlation helps: for Laplace
innovations at $t=0.05$ the median error falls from roughly one half at
$\alpha=0.5$ to $0.27$ at $\alpha=0.85$ and $0.13$ at $\alpha=0.95$, so
informative neighbours reduce the weight the local non-Gaussian shape
carries, an interpretation rather than a measured mechanism. The estimator is
fixed by the model class; the size of its error is each law's own.
